\chapter{核主成分分析}
\label{chap:kpca}
\begin{quote}\small
\textit{本章摘要。}核主成分分析在核诱导的特征空间中做PCA。本章从特征空间协方差算子出发，推导中心化Gram矩阵的特征分解、单位方向的归一化和新样本得分，并说明原空间重建为什么需要额外假设。
\end{quote}

线性PCA选择原坐标的线性组合。若希望描述的变化与输入平方值有关，就需要先改变表示。例如，一维输入$x$的线性坐标仍只能保留其线性变化，平方特征则可以表达与幅值大小有关、与正负号无关的变化。这个例子说明表示选择的差异，一维到一维本身不构成降维。核主成分分析（Kernel Principal Component Analysis，kPCA）先在$\phi(\bm x)$所在的特征空间中定义方差，再利用核内积计算主成分\cite{scholkopf1998kpca}。它使用上一章的确定性PCA目标，不自动提供概率生成模型。MML第10.5节的高维计算和第10.8节的核化讨论提供衔接，所给核方法网页列出了kPCA的核展开形式\cite{deisenroth2020mml,zheng2021kernel}。

\section{核矩阵的特征中心化}
设训练输入为$\bm x_i\in\R^d$，$i=1,\ldots,n$，其中$n\ge2$。取实值对称半正定核$k$和相应特征映射$\phi:\R^d\to\mathcal H$，令$K\in\R^{n\times n}$满足$K_{ij}=k(\bm x_i,\bm x_j)$。PCA要度量围绕均值的变化，因此先定义训练特征均值与中心化特征：
\begin{equation}
 \bar\phi=\frac1n\sum_i\phi(\bm x_i),\qquad
 \psi_i=\phi(\bm x_i)-\bar\phi.
\end{equation}
均值可能位于无限维空间，但它与有限个特征的内积仍可用核计算。展开两中心化特征的内积：
\begin{equation}
 \langle\psi_i,\psi_j\rangle
 =K_{ij}-\frac1n\sum_lK_{il}-\frac1n\sum_lK_{lj}
 +\frac1{n^2}\sum_{l,m}K_{lm}.
\end{equation}
用$\bm1\in\R^n$表示全1向量，得到
\begin{equation}
 H=I_n-\frac1n\bm1\bm1^{\mathsf T},\qquad K_c=HKH.
 \label{eq:kpca-centered-gram}
\end{equation}
$H$对样本维度去均值，满足$H^2=H$。$K_c$是中心化特征的Gram矩阵，仍然半正定，且$K_c\bm1=0$，所以其秩至多为$n-1$。对原始输入中心化一般不能代替这个操作：非线性映射通常不满足$\phi(\bm x-\bar{\bm x})=\phi(\bm x)-\bar\phi$。

\section{由协方差算子推导核矩阵特征问题}
在特征空间$\mathcal H$定义经验协方差算子
\begin{equation}
 T\bm u=\frac1n\sum_i\psi_i\langle\psi_i,\bm u\rangle_{\mathcal H}.
 \label{eq:kpca-covariance-operator}
\end{equation}
它将一个方向$\bm u\in\mathcal H$映为另一个特征空间向量。单位方向的投影方差为$\langle\bm u,T\bm u\rangle$。正特征值方向满足$T\bm u=\eta\bm u$，$\eta>0$。由此$\bm u=\eta^{-1}T\bm u$必在$\psi_1,\ldots,\psi_n$张成的空间中。

定义线性映射$\Phi_c:\R^n\to\mathcal H$，$\Phi_c\bm v=\sum_i v_i\psi_i$，其中$\bm v\in\R^n$收集展开系数。其伴随映射$\Phi_c^*:\mathcal H\to\R^n$返回各内积，即$\Phi_c^*\bm u=(\langle\psi_1,\bm u\rangle,\ldots,\langle\psi_n,\bm u\rangle)^{\mathsf T}$。因此$K_c=\Phi_c^*\Phi_c$及$T=\Phi_c\Phi_c^*/n$，两个乘积分别作用于样本系数空间和特征空间。若
\begin{equation}
 K_c\bm v_j=\kappa_j\bm v_j,\qquad
 \|\bm v_j\|_2=1,\quad\kappa_j>0,
\end{equation}
则取
\begin{equation}
 \bm u_j=\frac{\Phi_c\bm v_j}{\sqrt{\kappa_j}},\qquad
 T\bm u_j=\frac{\kappa_j}{n}\bm u_j,
 \qquad \|\bm u_j\|_{\mathcal H}^2
 =\frac{\bm v_j^{\mathsf T}K_c\bm v_j}{\kappa_j}=1.
 \label{eq:kpca-normalized-direction}
\end{equation}
因此核矩阵特征值$\kappa_j$与协方差特征值$\eta_j$相差因子$n$。写公式时需确认特征分解对象。只把核矩阵特征向量当作投影方向而不除以$\sqrt{\kappa_j}$，会改变主成分尺度。

训练样本$i$的第$j$个得分为
\begin{equation}
 z_{ij}=\langle\bm u_j,\psi_i\rangle
 =\frac{(K_c\bm v_j)_i}{\sqrt{\kappa_j}}
 =\sqrt{\kappa_j}\,v_{ij}.
 \label{eq:kpca-training-score}
\end{equation}
这里$v_{ij}$为$\bm v_j$的第$i$个分量。由于正特征值特征向量与$\bm1$正交，得分均值为零，方差为$\kappa_j/n$。令$r_c=\operatorname{rank}(K_c)$，保留方向数须满足$1\le q\le r_c$。特征空间平均平方重建误差为$\sum_{j=q+1}^{r_c}\kappa_j/n$，即被舍弃方向的方差之和。若$r_c=0$，训练特征没有可提取的非零方差方向。

\section{新样本的中心化与投影}
对新输入$\bm x$，定义$\bm k_{\bm x}\in\R^n$，其第$i$项为$k(\bm x_i,\bm x)$。新特征仍减去训练均值$\bar\phi$。中心化核向量是
\begin{equation}
 \tilde{\bm k}_{\bm x}=\bm k_{\bm x}-\frac1nK\bm1
 -\frac1n\bm1\bm1^{\mathsf T}\bm k_{\bm x}
 +\frac1{n^2}\bm1\bm1^{\mathsf T}K\bm1.
 \label{eq:kpca-test-centering}
\end{equation}
四项依次来自原始内积、训练特征均值、新样本与训练均值的内积，以及均值自身的平方长度。于是
\begin{equation}
 z_j(\bm x)=\frac{\bm v_j^{\mathsf T}\tilde{\bm k}_{\bm x}}{\sqrt{\kappa_j}}.
 \label{eq:kpca-new-score}
\end{equation}
不能把测试批次加入训练均值，否则同一查询点会因与它一起到达的测试样本不同而得到不同表示。交叉验证时，应在每个训练折内单独计算$K$、中心化统计量与特征分解。

\section{平方特征的完整算例}
取$x_1=-1,x_2=0,x_3=1$，核为$k(x,z)=x^2z^2$。它对应$\phi(x)=x^2$，是有效的半正定核。特征均值为$2/3$，中心化特征向量为$\bm r=(1/3,-2/3,1/3)^{\mathsf T}$。于是
\begin{equation}
 K_c=\bm r\bm r^{\mathsf T},\qquad
 \kappa_1=\|\bm r\|^2=2/3,\qquad
 \bm v_1=(1,-2,1)^{\mathsf T}/\sqrt6.
\end{equation}
式\eqref{eq:kpca-training-score}给出得分$(1/3,-2/3,1/3)$，经验方差为$2/9$。取新样本$x=2$，中心化特征为$4-2/3=10/3$，故$\tilde{\bm k}_{2}=(10/3)\bm r$。代入式\eqref{eq:kpca-new-score}得到$z_1(2)=10/3$。

图\ref{fig:kpca-square-map}展示这个映射。它合并了$x=-1$与$x=1$，保留的是大小相关的平方变化。若任务需要区分符号，这个核会丢失必要信息。
\begin{figure}[htbp]
\centering
\begin{tikzpicture}[node distance=0.65cm and 2.1cm]
 \node[idi input] (a) {$x=-1$};
 \node[idi input,below=of a] (b) {$x=0$};
 \node[idi input,below=of b] (c) {$x=1$};
 \node[idi state,right=of a] (u) {$\psi=1/3$};
 \node[idi state,right=of b] (v) {$\psi=-2/3$};
 \node[idi state,right=of c] (w) {$\psi=1/3$};
 \draw[idi arrow] (a)--(u); \draw[idi arrow] (b)--(v); \draw[idi arrow] (c)--(w);
 \node[right=0.6cm of v,align=left,font=\small] {映射：$\phi(x)=x^2$\\中心化：减去$2/3$};
\end{tikzpicture}
\caption{平方核的特征中心化。箭头表示具体计算，正负输入可能映为同一特征，核化不保证保留原输入的所有区别。}
\label{fig:kpca-square-map}
\end{figure}

\section{原像重建与核参数选择}
在特征空间中，重建可写成
\begin{equation}
 \hat\phi(\bm x)=\bar\phi+\sum_{j=1}^qz_j(\bm x)\bm u_j.
\end{equation}
这个向量不一定属于映射$\phi$的像，即可能没有任何原输入映到它。即使存在，也可能不唯一。平方特征例子中，特征值1同时对应输入$-1$和1。寻找原输入称为原像问题（pre-image problem），可考虑最小化$\|\phi(\bm x)-\hat\phi\|_{\mathcal H}^2$，或者另训练解码器。两者都增加了算法与假设，不能把核得分直接当成已恢复的原坐标。

使用高斯核时，$\ell$改变特征空间的几何。很小的$\ell$使不同训练点近似正交，中心化核接近$H$，多个特征值接近，主方向不稳定。很大的$\ell$使$K$接近全1矩阵，中心化后的变化趋小。应以独立验证的任务表现选择核和$q$。高核空间解释方差不保证低原空间重建误差。

完整核矩阵需要$O(n^2)$存储，完整特征分解常需$O(n^3)$计算。迭代特征求解器可只取前$q$个方向，低秩近似可进一步减少成本。特征值接近零时应设数值阈值避免除以很小的$\sqrt{\kappa_j}$，并把舍弃阈值与浮点误差联系起来。

\section{阅读依据}
本章结合MML第10.5与10.8节、核方法网页的核PCA部分及核PCA原始论文\cite{scholkopf1998kpca}。中心化、新样本变换和归一化在正文中完整展开。
